\section{Khảo sát và phân tích đề tài}

\subsection{Khảo sát thực trạng}

\subsubsection{Khảo sát quy trình giảng viên}

Nhóm tiến hành phỏng vấn và khảo sát online với \ldots\ giảng viên đại học,
tập trung vào ba câu hỏi chính: thời gian soạn nội dung kiểm tra, công cụ
đang sử dụng, và mức độ sẵn sàng dùng AI.

\begin{figure}[H]
    \centering
    % \includegraphics[width=0.9\textwidth]{figures/survey_instructor.png}
    \caption{Kết quả khảo sát giảng viên về quy trình tạo nội dung kiểm tra.}
    \label{fig:survey-instructor}
\end{figure}

\textbf{Nhận xét:} Quy trình thủ công hiện tại tiêu tốn trung bình
\ldots\ giờ/chương và hoàn toàn không có sự hỗ trợ của công cụ AI.
Đây là pain point rõ ràng mà hệ thống đề xuất cần giải quyết.

\subsubsection{Khảo sát học viên}

\begin{figure}[H]
    \centering
    % \includegraphics[width=0.9\textwidth]{figures/survey_learner.png}
    \caption{Kết quả khảo sát học viên về thói quen học tập.}
    \label{fig:survey-learner}
\end{figure}

\textbf{Nhận xét:} Phần lớn học viên mong muốn nội dung ngắn gọn và kiểm tra
tức thì --- phù hợp hoàn toàn với triết lý Micro-learning của hệ thống đề xuất.

\subsection{Khảo sát giải pháp tương tự}

\subsubsection{Coursera --- AI-generated Summaries}

Coursera gần đây bổ sung tính năng AI tóm tắt transcript bài giảng video,
hiển thị bên cạnh video. Tuy nhiên, tính năng này chỉ áp dụng cho video có
sẵn trên nền tảng và không cho phép giảng viên bên ngoài tải tài liệu lên để
xử lý.

\subsubsection{Khan Academy --- Khanmigo (AI Tutor)}

Khanmigo là chatbot AI tích hợp vào Khan Academy, hỗ trợ học viên giải thích
bài tập theo phong cách Socratic (đặt câu hỏi dẫn dắt thay vì đưa đáp án
thẳng). Giải pháp này không tự động sinh Quiz từ tài liệu mới và chỉ hỗ trợ
tiếng Anh.

\subsubsection{Quizlet AI --- Smart Notes}

Quizlet cho phép người dùng nhập hoặc paste văn bản, sau đó AI tự động tạo
flashcard và câu hỏi trắc nghiệm. Ưu điểm là đơn giản, nhanh. Nhược điểm:
không xử lý được file PDF/DOCX phức tạp, không tích hợp vào quy trình giảng
dạy (không có flow: upload $\to$ review $\to$ publish $\to$ học viên học).

\subsubsection{So sánh và đánh giá các giải pháp}

\begin{table}[H]
    \centering
    \caption{So sánh các giải pháp tương tự với hệ thống đề xuất.}
    \label{tab:compare-solutions}
    \begin{tabular}{p{3cm}p{2cm}p{2cm}p{2cm}p{2cm}p{3.5cm}}
        \toprule
        \textbf{Tiêu chí} & \textbf{Coursera AI} & \textbf{Khanmigo} &
        \textbf{Quizlet AI} & \textbf{Notion AI} & \textbf{Hệ thống đề xuất} \\
        \midrule
        Upload tài liệu riêng  & Không & Không & Có (text) & Có & Có (PDF/DOCX/PPTX) \\
        Sinh Quiz tự động       & Không & Không & Có        & Không & Có \\
        Review trước publish    & ---   & ---   & Không     & Không & Có \\
        Hỗ trợ tiếng Việt       & Không & Không & Một phần  & Có    & Có \\
        Giao diện học viên      & Có    & Có    & Có        & Không & Có \\
        Chi phí                 & Cao   & Cao   & Miễn phí + Pro & Có phí & Tự host \\
        \bottomrule
    \end{tabular}
\end{table}

\noindent\textbf{Kết luận:} Không có giải pháp nào hiện tại đáp ứng đầy đủ
ba tiêu chí: (1) upload tài liệu riêng, (2) có quy trình review Human-in-the-Loop,
và (3) hỗ trợ tiếng Việt. Đây là khoảng trống mà đề tài hướng đến lấp đầy.

\subsection{User Stories}

\subsubsection{Giảng viên}

\begin{longtable}{|p{0.5cm}|p{6cm}|p{8cm}|}
    \caption{User Stories của Giảng viên.} \label{tab:us-instructor} \\
    \hline
    \textbf{\#} & \textbf{User Story} & \textbf{Tiêu chí chấp nhận} \\
    \hline
    \endfirsthead \hline
    \textbf{\#} & \textbf{User Story} & \textbf{Tiêu chí chấp nhận} \\
    \hline \endhead

    1 & Là giảng viên, tôi muốn upload file PDF/DOCX để AI tự động tạo nội
        dung Micro-learning, giúp tôi tiết kiệm thời gian soạn bài. &
        Hệ thống nhận file $\leq$50MB; thông báo lỗi rõ nếu sai định dạng;
        pipeline khởi động $\leq$5 giây sau upload. \\ \hline

    2 & Là giảng viên, tôi muốn theo dõi tiến trình xử lý AI để biết khi
        nào có thể review. &
        Hiển thị 4 bước xử lý; cập nhật tự động; không cần refresh. \\ \hline

    3 & Là giảng viên, tôi muốn chỉnh sửa nội dung AI trước khi publish để
        đảm bảo chính xác học thuật. &
        Chỉnh sửa inline trực tiếp; lưu tức thì; có lịch sử sửa đổi. \\ \hline

    4 & Là giảng viên, tôi muốn publish khóa học chỉ khi đã review xong để
        học viên không thấy nội dung chưa được kiểm tra. &
        Nút Publish chỉ active sau khi $\geq$1 section được review. \\ \hline

    5 & Là giảng viên, tôi muốn xem thống kê học viên (ai đã học, điểm Quiz)
        để đánh giá hiệu quả nội dung. &
        Dashboard hiển thị: số học viên, tỷ lệ hoàn thành, điểm Quiz TB. \\ \hline
\end{longtable}

\subsubsection{Học viên}

\begin{longtable}{|p{0.5cm}|p{6cm}|p{8cm}|}
    \caption{User Stories của Học viên.} \label{tab:us-learner} \\
    \hline
    \textbf{\#} & \textbf{User Story} & \textbf{Tiêu chí chấp nhận} \\
    \hline
    \endfirsthead \hline
    \textbf{\#} & \textbf{User Story} & \textbf{Tiêu chí chấp nhận} \\
    \hline \endhead

    1 & Là học viên, tôi muốn xem danh sách khóa học và tiến độ cá nhân để
        lên kế hoạch học tập. &
        Progress bar cá nhân trên mỗi card khóa học; lọc theo trạng thái. \\ \hline

    2 & Là học viên, tôi muốn học nội dung dạng Card ngắn gọn để có thể
        học mọi lúc, mọi nơi. &
        Card $\leq$200 từ; swipe gesture (mobile); phím mũi tên (desktop). \\ \hline

    3 & Là học viên, tôi muốn làm Quiz và nhận phản hồi tức thì để biết
        mình đã hiểu đúng chưa. &
        Hiển thị đúng/sai + giải thích ngay sau khi chọn; không thể đổi đáp án. \\ \hline

    4 & Là học viên, tôi muốn tìm kiếm nội dung bằng câu hỏi tự nhiên
        để ôn tập đúng phần cần thiết. &
        Nhập query; hệ thống trả về top-5 Card liên quan trong $\leq$2 giây. \\ \hline
\end{longtable}

\subsubsection{Quản trị viên}

\begin{longtable}{|p{0.5cm}|p{6cm}|p{8cm}|}
    \caption{User Stories của Quản trị viên.} \label{tab:us-admin} \\
    \hline
    \textbf{\#} & \textbf{User Story} & \textbf{Tiêu chí chấp nhận} \\
    \hline
    \endfirsthead \hline
    \textbf{\#} & \textbf{User Story} & \textbf{Tiêu chí chấp nhận} \\
    \hline \endhead

    1 & Là quản trị viên, tôi muốn quản lý tài khoản người dùng (tạo, khóa,
        đổi vai trò) để vận hành hệ thống. &
        CRUD tài khoản; đổi vai trò INSTRUCTOR/LEARNER; khóa tài khoản vi phạm. \\ \hline

    2 & Là quản trị viên, tôi muốn xem log chi phí API để kiểm soát ngân
        sách sử dụng LLM. &
        Dashboard chi phí theo ngày/tháng; cảnh báo khi vượt ngưỡng. \\ \hline
\end{longtable}

\subsection{Yêu cầu nghiệp vụ}

Yêu cầu nghiệp vụ (Business Requirements --- BR) mô tả \textit{what} hệ
thống phải làm từ góc nhìn nghiệp vụ, không đề cập chi tiết kỹ thuật.

\subsubsection{BR-01: Xác thực và phân quyền người dùng}

Hệ thống phải phân biệt rõ ràng ba vai trò: Giảng viên, Học viên và Quản
trị viên. Mỗi vai trò chỉ được thực hiện các thao tác thuộc phạm vi của mình.
Người dùng phải xác thực trước khi truy cập bất kỳ chức năng nào.

\subsubsection{BR-02: Quản lý tài liệu và nội dung học tập}

Giảng viên phải có khả năng đưa tài liệu vào hệ thống và nhận lại nội dung
Micro-learning. Nội dung AI sinh ra phải có thể được giảng viên kiểm duyệt
và chỉnh sửa trước khi công bố cho học viên. Học viên chỉ thấy nội dung đã
được giảng viên chấp thuận.

\subsubsection{BR-03: Xử lý và sinh nội dung AI}

Hệ thống phải tự động phân tích cấu trúc tài liệu, chia thành các đơn vị
học tập có nghĩa, và sinh ra nội dung tóm tắt cùng câu hỏi kiểm tra phù hợp
với nội dung gốc. Nội dung sinh ra phải ở ngôn ngữ tương ứng với tài liệu
đầu vào.

\subsubsection{BR-04: Trải nghiệm học tập và đánh giá}

Học viên phải có thể học theo đơn vị nhỏ và tự đánh giá kiến thức sau mỗi
đơn vị. Hệ thống phải lưu lại tiến độ và kết quả của học viên để phục vụ
việc theo dõi và cải thiện.

\subsection{Yêu cầu tính năng}

Yêu cầu tính năng (Functional Requirements --- FR) cụ thể hóa BR thành các
chức năng hệ thống có thể đo lường.

\subsubsection{FR-01: Đăng nhập / Đăng xuất / Đăng ký tài khoản}

\begin{itemize}
    \item Người dùng đăng ký bằng email + mật khẩu; xác nhận qua email.
    \item Đăng nhập nhận JWT access token (1 giờ) + refresh token (7 ngày).
    \item Mật khẩu hash bằng \texttt{bcrypt}; không lưu plaintext.
    \item Đăng xuất vô hiệu hóa refresh token phía server.
\end{itemize}

\subsubsection{FR-02: Quản lý tài liệu}

\begin{itemize}
    \item Upload file PDF/DOCX/PPTX $\leq$ 50MB; kiểm tra MIME type.
    \item Lưu file vào MinIO; tạo bản ghi \texttt{Document} với trạng thái
          \texttt{QUEUED}.
    \item Hiển thị danh sách tài liệu của giảng viên với trạng thái xử lý
          theo thời gian thực.
    \item Xóa tài liệu: xóa file MinIO + cascade xóa chunk, card, quiz liên quan.
\end{itemize}

\subsubsection{FR-03: Xử lý tài liệu tự động (AI Pipeline)}

\begin{itemize}
    \item Sau khi upload, hệ thống tự động enqueue job và AI Pipeline worker
          xử lý bất đồng bộ.
    \item Trạng thái cập nhật theo thời gian thực qua server streaming:
          QUEUED $\to$ PARSING $\to$ CHUNKING $\to$ GENERATING $\to$ DONE
          (hoặc ERROR).
    \item Retry tối đa 3 lần nếu LLM API lỗi; ghi log chi tiết mỗi lần retry.
    \item Pipeline có thể xử lý đồng thời tối đa 10 tài liệu.
\end{itemize}

\subsubsection{FR-04: Review và publish khóa học}

\begin{itemize}
    \item Giảng viên xem, sửa (inline edit) và xóa từng Card và câu Quiz.
    \item Thêm câu hỏi thủ công vào bất kỳ section nào.
    \item Sắp xếp lại thứ tự Card trong section bằng drag-and-drop.
    \item Publish khóa học: điều kiện $\geq$ 1 section có \texttt{is\_reviewed=true}.
    \item Unpublish: ẩn khóa học khỏi danh sách học viên (soft delete).
\end{itemize}

\subsubsection{FR-05: Trải nghiệm học tập}

\begin{itemize}
    \item Học viên xem Card theo thứ tự; hệ thống lưu tiến độ sau mỗi Card.
    \item Làm Quiz: hiển thị ngẫu nhiên từ pool câu hỏi; phản hồi tức thì.
    \item Tìm kiếm ngữ nghĩa: vector similarity top-K.
    \item Xem lịch sử học: tiến độ theo khóa học, điểm Quiz theo thời gian.
\end{itemize}

\subsection{Yêu cầu phi tính năng}

\begin{longtable}{|p{1.5cm}|p{4cm}|p{9cm}|}
    \caption{Yêu cầu phi tính năng (NFR).} \label{tab:nfr} \\
    \hline
    \textbf{ID} & \textbf{Tên} & \textbf{Tiêu chí cụ thể} \\
    \hline
    \endfirsthead \hline
    \textbf{ID} & \textbf{Tên} & \textbf{Tiêu chí cụ thể} \\
    \hline \endhead

    NFR-01 & Hiệu năng &
        Xử lý tài liệu 20 trang $\leq$ 3 phút; API read $\leq$ 500ms ở P95;
        hỗ trợ 50 concurrent users xem Card/Quiz đồng thời. \\ \hline

    NFR-02 & Bảo mật &
        JWT trên mọi API; HTTPS bắt buộc; API key LLM qua biến môi trường;
        rate limiting 100 req/phút/IP; OWASP Top 10 check. \\ \hline

    NFR-03 & Độ tin cậy &
        Uptime $\geq$ 95\% staging; pipeline idempotent (retry an toàn);
        không mất dữ liệu khi lỗi giữa chừng. \\ \hline

    NFR-04 & Khả năng bảo trì &
        Kiến trúc DDD; unit test coverage $\geq$ 70\%; tài liệu API đầy đủ
        (auto-generated từ \texttt{.proto}). \\ \hline

    NFR-05 & Khả năng mở rộng &
        Thêm parser mới chỉ cần tạo adapter; thay LLM provider không sửa
        use case; Docker Compose $\to$ Kubernetes khi cần. \\ \hline

    NFR-06 & Khả dụng (Usability) &
        Không cần đào tạo; giao diện responsive; thông báo lỗi rõ và có
        hướng dẫn khắc phục; SUS Score $\geq$ 68/100. \\ \hline
\end{longtable}

\subsection{Tổng kết Chương 2}

Chương 2 đã trình bày toàn diện bức tranh hiện tại của bài toán: từ kết quả
khảo sát thực tế với giảng viên và học viên, phân tích so sánh bốn giải pháp
tương tự, đến đặc tả chi tiết user stories và các yêu cầu hệ thống (BR, FR,
NFR). Các phân tích này là căn cứ để thiết kế hệ thống trong các chương
tiếp theo.
